Time-series foundation models promise that you can pretrain once and then forecast any series without fitting.
The evidence comes almost entirely from mid-latitude, data-rich places, and often from reanalysis rather than
from instruments. BanglaWeatherBench asks what that promise is worth in a tropical monsoon climate, measured
against what stations actually recorded. It scores 8 zero-shot foundation models against 10 trained and
statistical baselines on 16 daily tasks over 170{,}571 rolling-origin windows, on 35 Bangladesh
Meteorological Department stations. Three controls make the result readable: 32 European stations from the same
network; NASA POWER reanalysis at the same 77 coordinates; and a context-gap control applied to models
and baselines alike. We resample at the forecast origin, because overlapping windows make ordinary tests
overconfident. The result is mixed, and the nulls are the useful part. Foundation models are the strongest family,
yet the best one is statistically indistinguishable from the best baseline in 29 of 48 comparisons. On station
observations it wins 3 of 24; on reanalysis at those coordinates, 11 of the same 24. Against a day-of-year
climatology on tropical rainfall nothing wins at all: not beyond a one-day lead at the stations, nor at any horizon
on ten-day satellite rainfall. Reanalysis looks more predictable than the stations it stands in for in 78
of 108 comparisons, so benchmarks built on it flatter models. At a 30-day lead 4 of 8 models lose
significantly more temperature skill in Bangladesh than a gap-matched control; roughly \SI{60}{\percent} of that
gap is sparser station context, not climate. Skill also reverses inside the monsoon: strong in break spells, worse than
climatology in active ones, when forecasts matter most. The whole leaderboard reproduces in 4.16 GPU-hours on a free
Tesla T4. Compute is not the barrier.
